\clearpage
\VerdierPhysicalPageAnchor{0142}%
\begingroup
\small
\setlength{\parskip}{0pt}%
\setlength{\emergencystretch}{1.5em}%

\VerdierSemanticAnchor{ch2-2-2-11}
\noindent\textbf{系 2.2.11.}\quad
\begingroup\itshape
\VerdierSemanticAnchor{ch2-2-2-11-setup}%
$\mathcal D$ を三角圏、$\mathcal B\subset\mathcal D$ を充満三角部分圏
とする。

\VerdierSemanticAnchor{ch2-2-2-11-a}
\par\noindent a) 商関手
$Q\colon\mathcal D\longrightarrow\mathcal D/\mathcal B$ の核は、
$\mathcal B$ を含む最小の厳密充満飽和三角部分圏である。

\VerdierSemanticAnchor{ch2-2-2-11-b-order-isomorphism}
\par\noindent b) 各厳密充満飽和三角部分圏にその飽和乗法的系 (2.1.8)
を対応させる写像 $\mathcal B\longmapsto S(\mathcal B)$ は、包含関係で
順序づけた厳密充満飽和三角部分圏の集合から、包含関係で順序づけた
$\mathcal D$ の飽和乗法的系の集合への順序集合の同型である。

\VerdierSemanticAnchor{ch2-2-2-11-b-inverse}
逆写像は、各飽和乗法的系に対応する局所化関手の核を対応させる。

\VerdierSemanticAnchor{ch2-2-2-11-c-factorization}
\par\noindent c) 核が部分圏 $\mathcal B$ を含む完全（それぞれ
コホモロジー的）関手 $F\colon\mathcal D\longrightarrow\mathcal D'$ は、
$Q\colon\mathcal D\longrightarrow\mathcal D/\mathcal B$ を通じて完全
（それぞれコホモロジー的）関手
$G\colon\mathcal D/\mathcal B\longrightarrow\mathcal D'$ として一意に
分解する。

\VerdierSemanticAnchor{ch2-2-2-11-c-morphism-lifting}
核が部分圏 $\mathcal B$ を含む完全（それぞれコホモロジー的）関手の射は、
それらの分解の間の射から一意に生じる。
\endgroup

\VerdierSemanticAnchor{ch2-2-2-11-easy-consequences}
これらは定理 (2.2.6) から容易に従う。

\VerdierSemanticAnchor{ch2-2-2-12}
\smallskip
\noindent\textbf{2.2.12. 性質 TRIV に関する注意。}
\VerdierSemanticAnchor{ch2-2-2-12-pretriangulated-category}%
性質 TRI、TRII、TRIII を満たす三角の族を備えた厳密な
$\mathbb Z$-圏を\emph{前三角圏}と呼ぶことにする。

\VerdierSemanticAnchor{ch2-2-2-12-induced-notions}
\emph{完全関手}、\emph{コホモロジー的関手}、\emph{乗法的系}の概念を、
三角圏の場合と同様に定義する。

\VerdierSemanticAnchor{ch2-2-2-12-localization}
$\mathcal D$ を前三角圏、$S$ をその乗法的系とする (2.1.4)。局所化圏
$\mathcal D(S^{-1})$ 上には、局所化関手が完全となるような前三角圏の
構造がただ一つ存在する。

\VerdierSemanticAnchor{ch2-2-2-12-kernel-subcategory}
局所化関手により零化される $\mathcal D$ の対象からなる充満部分圏は、
充満前三角部分圏である。

\VerdierSemanticAnchor{ch2-2-2-12-triv-necessity}
実際には、逆に、前三角部分圏に適切な飽和乗法的系を対応させる方法を
知る必要がある。編者が知る方法は性質 TRIV を用いるものだけである。

\VerdierSemanticAnchor{ch2-2-2-12-triv-role}
したがって、この本質的にかなり技術的な性質は、これまでのところ主に
定理 (2.2.6) から系 (2.2.11) を導くために用いられた。

\VerdierSemanticAnchor{ch2-2-3}
\medskip
\noindent\textbf{2.3. 局所化関手の性質。}

\VerdierSemanticAnchor{ch2-2-3-1}
\smallskip
\noindent\textbf{命題 2.3.1.}\quad
\itshape $\mathcal D$ を三角圏、$\mathcal B$ を充満三角部分圏とする。\upshape

\VerdierSemanticAnchor{ch2-2-3-1-p143-boundary}
\endgroup
